\section{Coupling the clocks to the scheduler}\label{sec:coupling}

An episode has two time coordinates. The physical time \(x\) locates the
anchors and probes of the scheduler. Within one bucket the difference
\(w\) of two clock readings will serve as time for an inner crossing.
We build a single blocking system, the \emph{driver}, in which that
inner crossing can be read from either clock vector or from the red-row
slacks. The readings themselves will be determined by the clock vectors.
Eliminating the readings preserves the rows tight on a path, the data
needed to transfer a crossing in Section~\ref{sec:transfer}.

\subsection{Physical time and reading time}

Fix a signature \(T\) and the signature \(T^+\) of
Theorem~\ref{thm:scheduler}. The two inputs are blocking systems
\(Q_{\mathrm{clock}}\) on a program on \(T\), and \(Q_{\mathrm{res}}\)
on a program \(\pi\) on \(T^+\) with \(m\) requests. In each input,
every viable set has more than \(J\) usable coordinates and every
admissible crossing of \([0,1]\) has at least \(D_*\) segments.
Choose a positive integer \(q_*\). Use Theorem~\ref{thm:scheduler}
to schedule \(\pi\) in \(N\) episodes in \([0,1]\) with \(B>4Vm\)
buckets. In the scheduler's bank of signature \(T\), place two fresh
copies of the clock program for each bucket. Write \(A,C\) for their
coordinate vectors in one bucket. Their coordinates are distinct
across buckets and from those of the scheduler.

For this bucket we give \(A\) and \(C\) lower bounds \(a(x)\) and
\(c(x)\). Their difference prescribes the reading when the clocks attain
these bounds. To construct it, shrink the closed probe bands within
their episodes, keeping each probe in the interior. In each episode
choose two distinct times in the interior of the bucket's probe band.
Prescribe a continuous piecewise affine function that is zero up to
the first time, increases linearly to one at the second time, and
remains one until the end of the episode. Call the closed intervals
where it is zero and one the \emph{early} and \emph{late} phases,
respectively. Mark an interval of positive length strictly inside the
intersection of the probe band with each of the phases. These two
\emph{overlaps} will host further crossings. Between episodes return
the prescribed function linearly to zero.

Write this piecewise affine function as \(a-c\), with both functions
convex. Give \(a\) the positive slope jumps of the function and \(c\)
the absolute values of its negative slope jumps, choosing initial
affine pieces with the prescribed difference. Each function has
\(O(N)\) pieces. Extend the end pieces affinely and add a common constant
so that \(a,c>1\) throughout \([0,1]\). Impose the coordinate bounds
\begin{equation}\label{eq:driver-floors}
                    A\ge a(x)\mathbf1,\qquad C\ge c(x)\mathbf1 .
\end{equation}
Choose a cap \(T_*\) strictly above all the floors of all buckets on
this interval.

On the reading axis \(w\) choose \(q_*\) successive triples of
disjoint, positive-length closed intervals in \((1/4,3/4)\). In each
triple name the intervals \(A,C,R\) in order. Call them \emph{work
windows}. Give each an open \emph{interval of influence} containing
it, with all closures pairwise disjoint. We build three blocking
systems on this common reading axis. The \emph{overlays} \(Q_A(w)\)
and \(Q_C(w)\) reproduce the clock input on their respective windows,
while \(Q_R(w)\) reproduces the resource input on the \(R\) windows.
Their coordinates will be supplied by the vectors in
\eqref{eq:driver-memberships} below.

\begin{lemma}[Copies on disjoint windows]\label{lem:overlay}
Each overlay uses $q_*$ affinely reparametrized and translated copies of
its input row list and $O(q_*)$ additional individual lower bounds per
coordinate, with the original coordinate set. On each window of its type
it coincides with the corresponding input after the increasing affine
change from the window to $[0,1]$ and translation by an affine function
of $w$ times $\mathbf1$. Every other row is strict there. For
$|w|\le 2T_*$, an input-copy row can be tight at a feasible point only
inside its own interval of influence.
\end{lemma}

\begin{proof}
Normalize the rows of the input as \(b\cdot y\ge\ell(w)\) with
\(b\ge0\) and \(b\cdot\mathbf1=1\). For the selected type of window,
order the windows and the intervening and exterior gaps along the line.
Assign an affine function to each window and gap as follows. Take their
slopes strictly increasing in this order. Choose successive intersection
times just outside the ends of the windows, inside the corresponding
intervals of influence, and determine the intercepts successively.
Their upper envelope visits the lines in the prescribed order. In
particular the line of a window is strictly above all the others
throughout that window, and strictly below the envelope outside its
interval of influence.

For a window with line \(p\) and affine change of time \(\phi\) onto
\([0,1]\), insert the entire input list with right sides
\(\ell(\phi(w))+p(w)\). For a gap with line \(p\), insert the individual
bounds \(y_i\ge p(w)\). We call these rows \emph{blanks}. On
\([-2T_*,2T_*]\) take \(M\ge1\) bounding all
\(\lvert\ell(\phi(w))\rvert\). Multiply the lines, but not the input
right sides \(\ell(\phi(w))\), by a common positive factor. By
compactness the strict gaps identified above can all be made larger
than \(2M\). On each piece of the envelope, individual rows of its own
copy or gap give \(y_i\ge p(w)-M\) for every \(i\). Since \(b\) is
a convex weight vector, all rows whose lines are lower by more than
\(2M\) are strict. On a work window the translated input alone
therefore implies every other row. Outside a copy's interval of
influence all its rows are strict at every feasible point.
\end{proof}

\subsection{Eliminating the clock readings}

The \(R\) overlay has one coordinate for each of the \(m\) requests in
\(\pi\). To supply that coordinate we use the smallest slack in its red
group. Keep the non-red scheduler rows. For every red row its slack is
the affine function left side minus right side, defined even when the
row is violated. For the current bucket \(P\) put
\[
       \mu_i=\min\{\text{slack of }h:h\text{ is a red row
       in the group }(P,i)\},\qquad 1\le i\le m.
\]
We replace the red rows by constraints joining these minima to the
two clock vectors. Choose \(\rho>0\) after fixing the overlays, and
then choose \(\tau>0\) after fixing \(\rho\). The smallness conditions
on these parameters are given in the proofs below.
The intended condition is the existence of $s,u\le T_*$ such that, for
$w=s-u$,
\begin{equation}\label{eq:driver-memberships}
\begin{aligned}
 z_A&=(A-s\mathbf1)/\rho                       &&\in Q_A(w),\\
 z_C&=(C-u\mathbf1)/\rho                       &&\in Q_C(w),\\
 z_R&=\bigl(\mu/\tau+(s+u)\mathbf1\bigr)/\rho    &&\in Q_R(w).
\end{aligned}
\end{equation}
Neither reading becomes a coordinate of the driver. To eliminate them,
write the $A$ and $C$ inequalities, respectively, as
\[
     s+\rho\ell_A(w)\le b_A\cdot A,\qquad
     u+\rho\ell_C(w)\le b_C\cdot C
\]
and count $s\le T_*$ and $u\le T_*$ as further upper inequalities of
their respective types. Each resource inequality has the form
\[
       b_R\cdot\mu/\tau+s+u\ge\rho\ell_R(w).
\]
The resource slack increases with the sum of the readings, while the
two upper slacks bound one reading each. This permits elimination with
positive multipliers, so tightness in the resulting row will record
tightness of all three constituents.

For a triple consisting of an upper inequality of each type and a
resource inequality, take their slacks with the feasible sign positive.
Add the three slacks, multiplying the two upper slacks by the unique
multipliers that cancel the coefficients of $s,u$ and leaving the
resource multiplier equal to one. At $\rho=0$ the two upper normals in
$(s,u)$ are $(-1,0),(0,-1)$ and the resource normal is $(1,1)$.
Thus both multipliers exist and are positive for sufficiently small
$\rho$, uniformly over the finite list of triples. For each
$i\in\operatorname{supp} b_R$ replace $\mu_i$ independently by the
slack of one row in its group and impose that the resulting sum be
nonnegative. This is one \emph{coupling row}. Its coordinate normal is
nonnegative and nonzero because it includes a positive multiple of at
least one red-row normal. Its right-hand side is affine in physical time.
The choices use at most $(2V)^{|\operatorname{supp} b_R|}$ rows for this
triple.

\begin{lemma}[Canonical readings and tight triples]\label{lem:coupling}
At every point satisfying the clock floors with $0\le x\le1$ there is a
canonical continuous choice of $(s,u)\in[0,T_*]^2$, determined by $A,C$,
satisfying all upper inequalities with at least one of each type tight.
It obeys
\begin{equation}\label{eq:canonical-estimate}
 \begin{aligned}
 |s-\min(T_*,\min_i A_i)|&\le C_*\rho,\\
 |u-\min(T_*,\min_i C_i)|&\le C_*\rho,
 \end{aligned}
 \qquad C_*\rho<1/32 .
\end{equation}
The coupling rows of the bucket are equivalent to
\eqref{eq:driver-memberships} at the canonical readings. They are also
equivalent to the existence of readings \(s,u\le T_*\) satisfying
\eqref{eq:driver-memberships}. At a point satisfying these equivalent
conditions, a coupling row is tight precisely when all three constituents
are tight there and every slack chosen in its expansion attains the
corresponding minimum. Deleting all coupling rows containing input-copy
rows from two or three different overlays preserves every slice with
$x\in[0,1]$ and the tight rows at every feasible point.
\end{lemma}

\begin{proof}
For a proposed pair $(s,u)$, evaluate $w=s-u$ and form the pair of bounds
whose first entry is
\[
      \min\bigl(T_*,\min_{(b,\ell)\text{ in }Q_A}
                              (b\cdot A-\rho\ell(w))\bigr)
\]
and whose second entry is given by $C,Q_C$ in place of $A,Q_A$. This map
is Lipschitz in the maximum norm with constant at most $2\rho$ times the
largest absolute slope of an upper row. On the square $[0,T_*]^2$ the
clock floors give $b\cdot A,b\cdot C>1$, and $\rho|\ell(w)|<1$ there for
small $\rho$. Hence the square maps to itself. The contraction has a unique fixed point. Its dependence on \(A,C\)
is continuous. Changing the clock vectors by at most \(\Delta\) in
maximum norm changes the map by at most \(\Delta\), hence changes the
fixed point by at most \(\Delta/(1-\kappa)\), where \(\kappa<1\)
is the Lipschitz bound. Each component of the fixed point attains its
minimum, giving one tight upper inequality of each type. Suppressing the terms $\rho\ell(w)$ changes the two
minima by at most $C_*\rho$, where $C_*$ bounds $|\ell|$ for
$|w|\le 2T_*$. The unperturbed minimum over $b\cdot A$ is $\min_i A_i$
because the normals are convex weight vectors and include all individual
coordinate normals. This proves \eqref{eq:canonical-estimate} after a
further decrease of $\rho$.

For any pair of readings, let $U_A,U_C,V_R$ be the three slacks of a
triple, with $V_R$ evaluated on the minima $\mu$. If its expansion
chooses red slacks $h_i$, the resulting physical row has slack
\begin{equation}\label{eq:coupling-slack}
  \lambda_A U_A+\lambda_C U_C+V_R+
    \frac{1}{\tau}\sum_{i\in\operatorname{supp} b_R}(b_R)_i(h_i-\mu_i),
\end{equation}
where $\lambda_A,\lambda_C>0$ are the elimination multipliers. The displayed
expression is independent of the readings. Any witness
to~\eqref{eq:driver-memberships} makes every summand nonnegative and so
gives all the expanded rows. Conversely, fix a resource row and use two
tight upper inequalities at the canonical readings. Choosing minimizing
red slacks in the expansion leaves only $V_R$ in~\eqref{eq:coupling-slack}.
Thus the expanded rows give every resource inequality at these readings.
At a feasible point all summands are nonnegative there, proving the
tightness criterion.

By Lemma~\ref{lem:overlay}, at the canonical $w\in[-T_*,T_*]$ two
different overlays cannot have tight input-copy rows simultaneously. Thus
each coupling row to be deleted is strict throughout the original slice.
For the deletion argument fix $x\in[0,1]$. Every row has a
nonzero nonnegative coordinate normal, so making all coordinates sufficiently
large gives a point strict in every original inequality. If a point
satisfied the retained rows and violated an original one, the segment
from the strict point would first meet a deleted equality at an original
feasible point, contradicting strictness there. The argument applies to
all buckets simultaneously. All tight expanded rows, in particular, have
been retained.
\end{proof}

\paragraph{Row count and activity.}
The driver consists of the non-red scheduler rows, the floors
\eqref{eq:driver-floors}, and the retained coupling rows for every bucket.
These inequalities define a blocking system on the whole real line,
with an individual bound on each coordinate. When this driver becomes an
input to a later overlay, each copy samples exactly $[0,1]$ on its work
window. Outside that window, the dominance proof applies directly to the
retained rows, whose affine right-hand sides are bounded on the compact
reading range. By Lemma~\ref{lem:coupling}, each surviving triple contains at most one
row from an input copy. With the numbers of blanks, caps and repetitions fixed, the number
of retained triples before expanding the minima is linear in the total
input row count. For each such triple the expansion costs at most
\((2V)^{|\operatorname{supp}b_R|}\) rows. Section~\ref{sec:recursion}
bounds this factor using the span of the resource row.

For a feasible driver point with $x\in[0,1]$, call a bucket
\emph{active} if one of its coupling rows is tight. By Lemma~\ref{lem:coupling} this is equivalent to a resource
inequality being tight at the canonical point. In that case every tight
upper or resource inequality occurs in a tight coupling row. Indeed, it
can be combined with tight rows of the other two types and minimizing slacks.
On the zero set of any coupling row in this feasible slab, its two upper
equalities determine $s,u$ as affine functions of
$A,C$ by the same nonsingular two-by-two system used in the elimination.

\begin{lemma}[Red slacks and active buckets]\label{lem:active}
Uniformly over feasible driver points with $0\le x\le1$, every red slack
is at least $-O(\tau)$, and every red slack occurring in a tight coupling
row is $O(\tau)$ in absolute value. For sufficiently small $\tau$, at
either anchor of any episode at most $2Vm$ buckets are active.
\end{lemma}

\begin{proof}
On the compact reading range, every coordinate of $z_R$ has a uniform
lower bound from its individual rows. In a tight resource inequality,
each coordinate with positive coefficient also has a uniform upper bound
because the right-hand side is bounded and all other terms are bounded below.
The lists of rows and coefficients are finite. Solving
\eqref{eq:driver-memberships} for $\mu$ gives $\mu_i\ge-O(\tau)$ for
every $i$ and $|\mu_i|=O(\tau)$ on each tight resource support. A red
slack used by a tight coupling row attains $\mu_i$. The constants can
depend on all the choices before $\tau$. Thus the scheduler coordinates
satisfy all its rows within its anchor tolerance, and every active
bucket has a red row within that tolerance of equality. The anchor
bound of Theorem~\ref{thm:scheduler} applies. Finitely many anchors and
red rows are involved, so a single choice of \(\tau\) works for all
of them.
\end{proof}
